\subsection{History and the remaining tensor problem}
The Hodge conjecture, formulated in Hodge's 1950 Congress address
\cite{Hodge1952}, asks which rational cohomology classes come from algebraic
cycles.  For a projective K3 surface itself, the Lefschetz $(1,1)$ theorem,
together with the classes of the surface and of a point, answers the question.
Products introduce a different problem: their Hodge classes include tensors
in the transcendental cohomologies of the factors.  Already for two factors,
the nontrivial part concerns Hodge homomorphisms between these cohomologies;
on higher powers, further invariant tensors occur.  We use the customary
background on K3 surfaces, their periods, and their transcendental Hodge
structures as presented in \cite{HuybrechtsK3}.

Kuga and Satake associated an abelian variety to a polarized Hodge structure
of K3 type \cite{KugaSatake}.  Its first cohomology carries a Clifford-algebra
representation, and the original weight-two Hodge structure embeds into a
tensor square of this weight-one structure.  Deligne developed this
construction in families and used it in his proof of the Weil conjecture
for K3 surfaces \cite[Sections~5--6]{DeligneK3}.  Andr\'e subsequently proved
that the Kuga--Satake correspondence is motivated
\cite[Section~7.1]{AndreMotives}.  These results provide the Hodge-theoretic
and motivic relation to abelian varieties; obtaining an algebraic cycle
representing the correspondence is the additional requirement in the
Kuga--Satake Hodge conjecture.  Van Geemen's account explains the Clifford
construction and its relation to exceptional Hodge classes on abelian
varieties \cite{VanGeemenKS}.

A second line of work constructs correspondences between K3 surfaces using
moduli spaces of sheaves.  Mukai's work supplied the fundamental examples
\cite{MukaiBundles}, and Nikulin extended the isometry result using
quadratic-lattice methods \cite{NikulinCorrespondences}.  Buskin proved that every rational Hodge isometry between
the second cohomologies of two projective K3 surfaces is algebraic
\cite[Theorem~1.1]{Buskin}; Huybrechts gave a proof using twisted derived
categories and established an isomorphism of their Chow motives
\cite[Theorems~0.1--0.2]{HuybrechtsIsogenous}.  The isometry hypothesis is
significant: arbitrary Hodge homomorphisms need not preserve the quadratic
forms.  Varesco proved that algebraic Kuga--Satake correspondences also make
Hodge similarities algebraic \cite[Theorem~4.5]{VarescoSimilarities}.

Zarhin determined the Hodge groups of transcendental K3 cohomology in terms
of its endomorphism field, which is either totally real or CM \cite{Zarhin}.
This description explains why a result about squares need not settle the
problem for powers.  In the totally real case, orthogonal invariant theory
produces determinant tensors in addition to the quadratic contractions.
Varesco identified the corresponding relative determinant classes and
proved a precise generation statement
\cite[Theorem~2.16]{VarescoPowers}.  When the endomorphism field is CM, the
invariants are generated in degree two; the work of Ram\'on Mar\'i, together
with algebraicity of Hodge isometries, gives the Hodge conjecture for all
powers \cite{RamonMari,Buskin}; see also
\cite[Theorem~2.2 and Corollary~2.3]{VarescoPowers}.  Here a CM endomorphism
field is the relevant condition; it need not make the whole transcendental
Hodge structure of CM type.

Algebraic Kuga--Satake correspondences have also been constructed for
important special families.  Paranjape treated the K3 surfaces arising
from double covers of the plane branched along six lines \cite{Paranjape},
and Schlickewei used their Kuga--Satake varieties to establish the Hodge
conjecture for their squares \cite[Theorem~2]{Schlickewei}.  Varesco proved an
all-powers theorem for four-dimensional families of generic Picard number
16 with a specified rational transcendental quadratic space, assuming
algebraicity of the Kuga--Satake correspondences throughout the family
\cite[Theorem~0.2]{VarescoPowers}.  The quadratic-space and family
hypotheses in that statement are essential to its scope.

Recent geometric constructions considerably enlarge these cases.
Floccari proved algebraicity of the Kuga--Satake correspondence and the
Hodge conjecture for all powers of the K3 surfaces associated with varieties
of generalized Kummer type \cite[Theorem~1.2]{FloccariKummer}.
Floccari and Fu obtained related results through singular varieties of
O'Grady's six-dimensional deformation type; their construction also gives
products of powers of the abelian, K3, and hyper-K\"ahler varieties
associated with one such singular variety
\cite[Theorem~5.7 and Corollary~5.8]{FloccariFu}.  These theorems illustrate
how geometric relations with abelian varieties control exceptional tensors
on particular families.  For a general list of K3 surfaces, a further
issue is to control the joint Hodge group: algebraicity of the tensors on
each individual factor does not by itself determine the mixed invariants.

